
\subsection{Static Analysis - Type System}
Let $\Gamma$ be the type environment set consisting of elements either $identifier: Type$ or $Type$. The typing rules have been organized in distinct categories for readability. The notation $Type[x, y, z]$ denotes refinements on $Type$ either through instantiating parametric (generic) types or restricting values by way of \texttt{trait} clauses.

Type refinements may be applicable to types with certain traits (ex., min/max values for ordered types, definedness of expressions with certain values, set sizes). % Is there a way to formalize this? Refinement calculus?

\subsubsection{Base Environment Interactions}
\paragraph{Includes:} Mechanisms for adding/extracting variables from the type environment.

\begin{gather}
\tag{Env $\emptyset$}
\begin{prooftree}
\infer0{\emptyset \vdash \diamond}
\end{prooftree}
\\
\tag{Env $I$}
\begin{prooftree}
\hypo{\Gamma \vdash A}
\hypo{I \not\in dom(\Gamma)}
\infer2{\Gamma, I:A \vdash \diamond}
\end{prooftree}
\\
\tag{Fetch Identifier}
\begin{prooftree}
\hypo{\Gamma_1, I:A, \Gamma_2 \vdash \diamond}
\infer1{\Gamma_1, I:A, \Gamma_2 \vdash I:A}
\end{prooftree}
\end{gather}

\subsubsection{Subtype Interactions}
\paragraph{Includes:} Generic subtyping rules.

\begin{gather}
\tag{Reflexive Subtype}
\begin{prooftree}
\hypo{\Gamma \vdash A}
\infer1{\Gamma \vdash A \leq A}
\end{prooftree}
\\
\tag{Transitive Subtype}
\begin{prooftree}
\hypo{\Gamma \vdash A \leq B}
\hypo{\Gamma \vdash B \leq C}
\infer2{\Gamma \vdash A \leq C}
\end{prooftree}
\\
\tag{Subsumption}
\begin{prooftree}
\hypo{\Gamma \vdash a:A}
\hypo{\Gamma \vdash A \leq B}
\infer2{\Gamma \vdash a: B}
\end{prooftree}
\\
\tag{Top Type}
\begin{prooftree}
\hypo{\Gamma \vdash \diamond}
\infer1{\Gamma \vdash Any}
\end{prooftree}
\\
\tag{Sub Top Type}
\begin{prooftree}
\hypo{\Gamma \vdash A}
\infer1{\Gamma \vdash A \leq Any}
\end{prooftree}
\\
\tag{Sub Function}
\begin{prooftree}
\hypo{\Gamma \vdash A' \leq A}
\hypo{\Gamma \vdash B \leq B'}
\infer2{\Gamma \vdash A \rightarrow B \leq A' \rightarrow B'}
\end{prooftree}
\\
\tag{Sub Set}
% TODO check this - doesnt match up with the definition of set as set: A -> bool, but does work intuitively
% ex. set[nat] <= set[int]
\begin{prooftree}
\hypo{\Gamma \vdash A \leq B}
\infer1{\Gamma \vdash set[A] \leq set[B]}
\end{prooftree}
\\
\tag{Sub Product}
\begin{prooftree}
\hypo{\Gamma \vdash A_1 \leq B_1}
\hypo{\Gamma \vdash A_2 \leq B_2}
\infer2{\Gamma \vdash A_1 \times A_2 \leq B_1 \times B_2 }
\end{prooftree}
\\
\tag{Sub Record}
\begin{prooftree}
\hypo{\Gamma \vdash A_1 \leq B_1 ... A_n \leq B_n}
\hypo{\Gamma \vdash A_{n+1} ... A_{n+m}}
\infer2{\Gamma \vdash  Record[I_1: A_1, ..., I_n: A_n, I_{n+1}: A_{n+1}, ... I_{n+m}: A_{n+m}] \leq  Record[I_1: B_1, ..., I_n: B_n]}
\end{prooftree}
\\
\tag{Type Refinement}
\begin{prooftree}
\hypo{\Gamma \vdash T}
\infer1{\Gamma \vdash T[\text{with } x] \leq T}
\end{prooftree}
\end{gather}

\subsubsection{The Set Type}
\paragraph{Includes:} Universal type; everything is a set. The empty set is a part of every type (aside from primitives).

\begin{gather}
\tag{Powerset}
\begin{prooftree}
\hypo{\Gamma \vdash T}
\infer1{\Gamma \vdash set[T]}
\end{prooftree}
\\
\tag{Emptyset}
\begin{prooftree}
\hypo{\Gamma \vdash \diamond}
\infer1{\Gamma \vdash EmptySet}
\end{prooftree}
\\
\tag{Emptyset Bottom}
\begin{prooftree}
\hypo{\Gamma \vdash EmptySet, set[T]}
\infer1{EmptySet \leq set[T]}
\end{prooftree}
\\
\tag{Set Enumeration}
\begin{prooftree}
\hypo{\Gamma \vdash e_1, ..., e_n: A}
\infer1{\Gamma \vdash \Set{e_1, ..., e_n}: set[A]}
\end{prooftree}
\end{gather}

\subsubsection{Primitive Types}
\paragraph{Includes:} Basic types built in to the language, like numbers (int and float), strings, and booleans. Although these could all be represented as sets, it is useful to distinguish these common types from other collections or user-defined types.

\begin{gather}
\tag{Primitives - bool}
\begin{prooftree}
\hypo{\Gamma \vdash \diamond}
\infer1{\Gamma \vdash bool }
\end{prooftree}
\\
\tag{Primitives - int}
\begin{prooftree}
\hypo{\Gamma \vdash \diamond}
\infer1{\Gamma \vdash int }
\end{prooftree}
\\
\tag{Primitives - float}
\begin{prooftree}
\hypo{\Gamma \vdash \diamond}
\infer1{\Gamma \vdash float }
\end{prooftree}
\\
\tag{Primitives - str}
\begin{prooftree}
\hypo{\Gamma \vdash \diamond}
\infer1{\Gamma \vdash str }
\end{prooftree}
\\
\tag{Nat from int}
\begin{prooftree}
\hypo{\Gamma \vdash int}
\infer1{\Gamma \vdash nat \leq int[\text{with } min=0] }
\end{prooftree}
\\
\tag{Int from float}
\begin{prooftree}
\hypo{\Gamma \vdash int, float}
\infer1{int \leq float }
\end{prooftree}
\end{gather}

\subsubsection{Collection Types}
\paragraph{Includes:} Sequences, relations, enums (frozen sets), and bags (multisets); Syntactic sugar for abstract data types defined in terms of sets.

\begin{gather}
\tag{Tuples}
\begin{prooftree}
\hypo{\Gamma \vdash T_1, ..., T_n}
\infer1{\Gamma \vdash tuple[T_1, ..., T_n] \quad tuple[T_1, ..., T_n] = T_1 \times ... \times T_n}
\end{prooftree}
\\
\tag{Empty Tuple}
\begin{prooftree}
\hypo{\Gamma \vdash \diamond}
\infer1{\Gamma \vdash (): EmptyTuple}
\end{prooftree}
\\
\tag{Empty Tuple Bottom}
\begin{prooftree}
\hypo{\Gamma \vdash EmptyTuple, T}
\infer1{\Gamma \vdash EmptyTuple \leq tuple[T]}
\end{prooftree}
\\
\tag{Tuple Enumeration}
\begin{prooftree}
\hypo{\Gamma \vdash e_1:A_1, ..., e_n: A_n}
\infer1{\Gamma \vdash ( e_1, ..., e_n ) : tuple[A_1, ..., A_n]}
\end{prooftree}
\\
\tag{Relation Type}
\begin{prooftree}
\hypo{\Gamma \vdash set[T_1], set[T_2]}
\infer1{\Gamma \vdash T_1 \rel T_2 \quad T_1 \rel T_2 = set[tuple[T_1, T_2]]}
\end{prooftree}
\\
\tag{Bag Type}
\begin{prooftree}
\hypo{\Gamma \vdash T, nat}
\infer1{\Gamma \vdash bag[T] \leq T \pfun nat }
\end{prooftree}
\\
\tag{Bag Enumeration}
\begin{prooftree}
\hypo{\Gamma \vdash e_1, ..., e_n: A}
\infer1{\Gamma \vdash \Bag{e_1, ..., e_n}: bag[A]}
\end{prooftree}
\\
\tag{Sequence Type}
\begin{prooftree}
\hypo{\Gamma \vdash T, nat}
\infer1{\Gamma \vdash sequence[T] \leq nat \pfun T }
\end{prooftree}
\\
\tag{Sequence Enumeration}
\begin{prooftree}
\hypo{\Gamma \vdash e_1, ..., e_n: A}
\infer1{\Gamma \vdash \langle e_1, ..., e_n \rangle : sequence[A]}
\end{prooftree}
\\
\tag{Enum from static set}
\begin{prooftree}
\hypo{\Gamma \vdash s_1, ..., s_n : str}
\hypo{A = \Set{s_1, ..., s_n} }
\hypo{immutable(A)}
\infer3{\Gamma \vdash enum[A] \leq set[A, \text{with } immutable]}
\end{prooftree}
\end{gather}



\subsubsection{Relation Subtypes}
\paragraph{Includes:} Relation subtypes examine four key properties of relations: totality ($t$), surjectivity ($t_r$), one-to-manyness ($1-$), and many-to-oneness ($1-_r$). Syntactic sugar to produce refined relations with these properties follows the definitions presented in Event-B. See Section 5.1.3 for more details.

\begin{gather}
\tag{Relation Subtype - Total Relation}
\begin{prooftree}
\hypo{\Gamma \vdash set[T_1], set[T_2]}
\infer1{\Gamma \vdash T_1 \trel T_2 \quad T_1 \trel T_2 = (T_1 \rel T_2)[\text{with } t]}
\end{prooftree}
\\
\tag{Relation Subtype - Surjective Relation}
\begin{prooftree}
\hypo{\Gamma \vdash set[T_1], set[T_2]}
\infer1{\Gamma \vdash T_1 \srel T_2 \quad T_1 \srel T_2 = (T_1 \rel T_2)[\text{with } t_r]}
\end{prooftree}
\\
\tag{Relation Subtype - Total Surjective Relation}
\begin{prooftree}
\hypo{\Gamma \vdash set[T_1], set[T_2]}
\infer1{\Gamma \vdash T_1 \strel T_2 \quad T_1 \strel T_2 = (T_1 \rel T_2)[\text{with } t, \text{ with } t_r]}
\end{prooftree}
\\
\tag{Relation Subtype - Partial Function}
\begin{prooftree}
\hypo{\Gamma \vdash set[T_1], set[T_2]}
\infer1{\Gamma \vdash T_1 \pfun T_2 \quad T_1 \pfun T_2 = (T_1 \rel T_2)[\text{with } 1-]}
\end{prooftree}
\\
\tag{Relation Subtype - Total Function}
\begin{prooftree}
\hypo{\Gamma \vdash set[T_1], set[T_2]}
\infer1{\Gamma \vdash T_1 \tfun T_2 \quad T_1 \tfun T_2 = (T_1 \rel T_2)[\text{with } t, \text{ with } 1-]}
\end{prooftree}
\\
\tag{Relation Subtype - Partial Injection}
\begin{prooftree}
\hypo{\Gamma \vdash set[T_1], set[T_2]}
\infer1{\Gamma \vdash T_1 \pinj T_2 \quad T_1 \pinj T_2 = (T_1 \rel T_2)[\text{with } 1-, \text{ with } 1-_r]}
\end{prooftree}
\\
\tag{Relation Subtype - Total Injection}
\begin{prooftree}
\hypo{\Gamma \vdash set[T_1], set[T_2]}
\infer1{\Gamma \vdash T_1 \tinj T_2 \quad T_1 \tinj T_2 = (T_1 \rel T_2)[\text{with } t, \text{ with } 1-, \text{ with } 1-_r]}
\end{prooftree}
\\
\tag{Relation Subtype - Partial Surjection}
\begin{prooftree}
\hypo{\Gamma \vdash set[T_1], set[T_2]}
\infer1{\Gamma \vdash T_1 \psur T_2 \quad T_1 \psur T_2 = (T_1 \rel T_2)[\text{with } t_r, \text{ with } 1-]}
\end{prooftree}
\\
\tag{Relation Subtype - Total Surjection}
\begin{prooftree}
\hypo{\Gamma \vdash set[T_1], set[T_2]}
\infer1{\Gamma \vdash T_1 \tsur T_2 \quad T_1 \tsur T_2 = (T_1 \rel T_2)[\text{with } t, \text{ with } t_r, \text{ with } 1-]}
\end{prooftree}
\\
\tag{Relation Subtype - Bijection}
\begin{prooftree}
\hypo{\Gamma \vdash set[T_1], set[T_2]}
\infer1{\Gamma \vdash T_1 \tbij T_2 \quad T_1 \tbij T_2 = (T_1 \rel T_2)[\text{with } t, \text{ with } t_r, \text{ with } 1-, \text{ with } 1-_r]}
\end{prooftree}
\end{gather}

\subsubsection{Simple Commands}
\paragraph{Includes:} Assignment, type assignment, type refinements, simple programming language commands/statements.

\begin{gather}
\tag{Variable Assignment}
\begin{prooftree}
\hypo{\Gamma \vdash E:A}
\hypo{I \not\in dom(\Gamma) \lor \Gamma \vdash I:A}
\infer2{\Gamma \vdash I:A := E }
\end{prooftree}
\\
\tag{Type Alias Assignment}
\begin{prooftree}
\hypo{\Gamma \vdash T}
\hypo{I \not\in dom(\Gamma)}
\infer2{\Gamma \vdash I := T }
\end{prooftree}
\\
\tag{Type Alias}
\begin{prooftree}
\hypo{\Gamma \vdash I := T}
\infer1{\Gamma, I \vdash \diamond \quad I = T}
\end{prooftree}
\\
\tag{Refined Variable Assignment}
\begin{prooftree}
\hypo{\Gamma \vdash E: A, E_1: A_1, ..., E_n:A_n}
\hypo{(I \not\in dom(\Gamma) \lor \Gamma \vdash I:A)}
\infer2{\Gamma \vdash I: A[\text{with } E_1: A_1, ..., \text{ with } E_n: A_n] := E}
\end{prooftree}
\\
\tag{Refined Type Alias}
\begin{prooftree}
\hypo{\Gamma \vdash T, E_1: A_1, ..., E_n: A_n}
\hypo{I \not\in dom(\Gamma)}
\infer2{\Gamma \vdash I := T[\text{with } E_1: A_1, ...,\text{ with } E_n: A_n]}
\end{prooftree}
\\
\tag{Command - break}
\begin{prooftree}
\hypo{\Gamma \vdash \diamond}
\infer1{\Gamma \vdash \text{break}: C}
\end{prooftree}
\\
\tag{Command - continue}
\begin{prooftree}
\hypo{\Gamma \vdash \diamond}
\infer1{\Gamma \vdash \text{continue}: C}
\end{prooftree}
\\
\tag{Command - skip}
\begin{prooftree}
\hypo{\Gamma \vdash \diamond}
\infer1{\Gamma \vdash \text{skip}: C}
\end{prooftree}
\\
\tag{Command - return}
\begin{prooftree}
\hypo{\Gamma \vdash E:A}
\infer1{\Gamma \vdash (\text{return } E): C}
\end{prooftree}
\end{gather}

\subsubsection{Quantification Expressions}
\paragraph{Includes:} Set, sequence, relation, and bag comprehensions, in addition to (set based) lambda expressions. Comprehension variables may be bound through either an $Identifier$ or an arbitrarily nested tuple of $Identifiers$

\begin{gather}
\tag{Lambda Expression}
\begin{prooftree}
\hypo{\Gamma, I_1:T_1, ..., I_n: T_n \vdash E:A, P: bool}
\hypo{I_i \notin dom(\Gamma)}
\hypo{i \neq j \implies I_i \neq I_j}
\infer3{\Gamma \vdash (\Lambda{I_1, ..., I_n}{P}{E}): (T_1, ..., T_n) \pfun A}
\end{prooftree}
\\
\tag{Quantifier}
\begin{prooftree}
\hypo{\Gamma \vdash G: Generator[I]}
\hypo{\Gamma , I \vdash E: A}
\infer2{\Gamma \vdash (G \mid E) : Quantifier[A]}
\end{prooftree}
\\
\tag{Branching Quantifier}
\begin{prooftree}
\hypo{\Gamma \vdash Q_i : Quantifier[A]}
\infer1{\Gamma \vdash (Q_1 ` ... ` Q_n): Quantifier[A]}
\end{prooftree}
\\
\tag{Generator nest}
\begin{prooftree}
\hypo{\Gamma \vdash G_i: Generator[T_i]}
\infer1{\Gamma \vdash G_1, ..., G_n : Generator[(T_1,..., T_n)]} % acts like a generator over the cartesian product
\end{prooftree}
\\
\tag{Generator}
\begin{prooftree}
\hypo{\Gamma \vdash S: set[T], P: bool}
\hypo{I = x_1:T_1, ..., x_n:T_n}
\hypo{\Gamma \vdash structuralMatch(I, T)}
\infer3{\Gamma \vdash (I \in S \cdot P): Generator[T]}
\end{prooftree}
\\
% \tag{Quantification Body \footnote{where $TupleIdentifier$ is a recursive tuple of either $Identifier$s or nested $TupleIdentifiers$, $structuralMatch(e,T)$ is a recursive match between identifiers in $e$ and the set type $T$, $hasGenerator(P,e \in g)$ means the expression $e \in g$ occurs in one of the top-level OR/AND clauses, and $structuralMatchInAllORs(e, T, P)$ means all top-level OR clauses must see $e \in g$ for any $g$ with type $T$. For now, we may just force the generator to appear in the first clause of the predicate, where the predicate is a top-level AND.}}
% \begin{prooftree}
% \hypo{I = x_1:T_1, ..., x_n:T_n}
% \hypo{\Gamma \vdash binds(P, I)}
% \hypo{\Gamma, I \vdash P:bool, E:A}
% \infer3{\Gamma \vdash (x_1, ..., x_n \cdot P \mid E): QuantificationBody[I, A]}
% \end{prooftree}
% \\
% \tag{Binds (with generator)}
% \begin{prooftree}
% \hypo{\Gamma \vdash g : set[T]}
% \hypo{\Gamma \vdash structuralMatch(I, T)}
% \infer2{\Gamma \vdash binds(I \in g, I)}
% \end{prooftree}
% \\
% \tag{Binds with OR}
% \begin{prooftree}
% \hypo{\Gamma \vdash binds(P_1, I), ..., binds(P_n, I)}
% \infer1{\Gamma \vdash binds(\bigvee P_i, I)}
% \end{prooftree}
% \\
% \tag{Binds with AND}
% \begin{prooftree}
% \hypo{\Gamma, I \vdash Q: bool}
% \hypo{\Gamma \vdash binds(P, I)}
% \infer2{\Gamma \vdash binds(P \land Q, I)}
% \end{prooftree}
% \\
\tag{Structural Match}
\begin{prooftree}
\hypo{\Gamma \vdash T}
\hypo{x \notin dom(\Gamma)}
\infer2{\Gamma \vdash structuralMatch(x, T)}
\end{prooftree}
\\
\tag{Structural Match with Tuple}
\begin{prooftree}
\hypo{\Gamma \vdash structuralMatch(x_1, T_1), ..., structuralMatch(x_n, T_n)}
\infer1{\Gamma \vdash structuralMatch((x_1, ..., x_n), (T_1, ..., T_n))}
\end{prooftree}
\\
\tag{Set Comprehension}
\begin{prooftree}
\hypo{\Gamma \vdash Q: Quantifier[A]}
\infer1{\Gamma \vdash \Set{Q}: set[A]}
\end{prooftree}
\\
\tag{Bag Comprehension}
\begin{prooftree}
\hypo{\Gamma \vdash Q: Quantifier[A]}
\infer1{\Gamma \vdash \Bag{Q}: bag[A]}
\end{prooftree}
\\
\tag{Sequence Comprehension}
\begin{prooftree}
\hypo{\Gamma \vdash Q: Quantifier[A]}
\infer1{\Gamma \vdash \List{Q}: sequence[A]}
\end{prooftree}
\\
\tag{Fold}
\begin{prooftree}
\hypo{\Gamma, i: A \vdash Q: Quantifier[A]}
\hypo{\Gamma \vdash Id: A}
\hypo{i \notin dom(\Gamma)}
\infer3{\Gamma \vdash (fold\ i := Id : Q): A}
\end{prooftree}
\\
\tag{Iter Quantifier}
\begin{prooftree}
\hypo{\Gamma \vdash G: IterGenerator[I], x_i:T_i, \text{block }C}
\infer1{\Gamma \vdash (G \rightarrow (x_1, ..., x_n) \mid C): IterQuantifier[(T_1, ..., T_n)]}
\end{prooftree}
\\
\tag{Iter Generator}
\begin{prooftree}
\hypo{\Gamma \vdash G: Generator[I], y_i:T_i}
\hypo{x_i \notin dom(\Gamma)}
\infer2{\Gamma \vdash (G : x_1 := y_1; ... ; x_n := y_n): IterGenerator[I]}
\end{prooftree}
\\
\tag{Iter}
\begin{prooftree}
\hypo{\Gamma \vdash Q: IterQuantifier[A]}
\infer1{\Gamma \vdash (iter\ Q): A}
\end{prooftree}
\end{gather}

Let $\otimes = \Set{ \bigcup, \bigcap }$

\begin{gather}
\tag{General Quantifiers}
\begin{prooftree}
\hypo{\Gamma \vdash Q: Quantifier[A]} % where A is a set
\infer1{\Gamma \vdash (\otimes Q): A}
\end{prooftree}
\end{gather}

Let $\otimes = \Set{ \forall, \exists }$

\begin{gather}
\tag{Bool Quantifiers}
\begin{prooftree}
\hypo{\Gamma \vdash Q: Quantifier[bool]}
\infer1{\Gamma \vdash (\otimes Q): bool}
\end{prooftree}
\end{gather}

Let $\otimes = \Set{ \sum, \prod }$

\begin{gather}
\tag{Numeric Quantifiers}
\begin{prooftree}
\hypo{\Gamma \vdash Q: Quantifier[T]}
\hypo{\Gamma \vdash T \in \Set{int, nat, float}}
\infer2{\Gamma \vdash (\otimes Q): T}
\end{prooftree}
\end{gather}

\subsubsection{Boolean Operations}
\paragraph{Includes:} Common predicate logic operations (equivalence, implication, and, or, etc.).

Let $\otimes = \Set{\equiv, \not\equiv, \implies, \land, \lor }$

\begin{gather}
\tag{Binary Boolean Operations}
\begin{prooftree}
\hypo{\Gamma \vdash b_1, b_2: bool}
\infer1{\Gamma \vdash b_1 \otimes b_2 : bool}
\end{prooftree}
\\
\tag{Boolean Operations - Negation}
\begin{prooftree}
\hypo{\Gamma \vdash b_1: bool}
\infer1{\Gamma \vdash \lnot b_1 : bool}
\end{prooftree}
\end{gather}

\subsubsection{Equality and Membership}
\paragraph{Includes:} Equality and ordering comparisons, set membership.

Let $\otimes = \Set{=, \neq}$

\begin{gather}
\tag{Equals}
\begin{prooftree}
\hypo{\Gamma \vdash a_1: T_1, a_2: T_2}
\infer1{\Gamma \vdash a_1 \otimes a_2 : bool}
\end{prooftree}
\end{gather}

Let $\otimes = \Set{<, \leq, >, \geq }$

\begin{gather}
\tag{Ordering Operators}
\begin{prooftree}
\hypo{\Gamma \vdash a_1, a_2: T}
\hypo{T \in \Set{int, float, nat}}
\infer2{\Gamma \vdash a_1 \otimes a_2 : bool}
\end{prooftree}
\end{gather}

Let $\otimes = \Set{\in, \not\in }$

\begin{gather}
\tag{Set Membership}
\begin{prooftree}
\hypo{\Gamma \vdash a_1: T_1, a_2: set[T_2]}
\infer1{\Gamma \vdash a_1 \otimes a_2: bool}
\end{prooftree}
\end{gather}

\subsubsection{Set Operations}
\paragraph{Includes:} Subset comparisons, union, intersection, cartesian product, etc. Also includes Maplet and Range.

Let $\otimes = \Set{\subset, \subseteq, \supset, \supseteq, \not\subset, \not\subseteq, \not\supset, \not\supseteq }$

\begin{gather}
\tag{Set Ordering Operations}
\begin{prooftree}
\hypo{\Gamma \vdash a_1: set[T_1], a_2: set[T_2]}
\infer1{\Gamma \vdash a_1 \otimes a_2: bool}
\end{prooftree}
\end{gather}

Let $\otimes = \Set{\cup, \cap, \setminus }$

\begin{gather}
\tag{Set Operations}
\begin{prooftree}
\hypo{\Gamma \vdash a_1, a_2: set[T]}
\infer1{\Gamma \vdash a_1 \otimes a_2: set[T]}
\end{prooftree}
\\
\tag{Cartesian Product}
\begin{prooftree}
\hypo{\Gamma \vdash a_1: set[T_1], a_2: set[T_2]}
\infer1{\Gamma \vdash a_1 \times a_2: T_1 \rel T_2}
\end{prooftree}
\\
\tag{Maplet}
\begin{prooftree}
\hypo{\Gamma \vdash a_1: T_1, a_2: T_2}
\infer1{\Gamma \vdash a_1 \mapsto a_2: tuple[T_1, T_2]}
\end{prooftree}
\\
\tag{Numerical Range}
\begin{prooftree}
\hypo{\Gamma \vdash a_1, a_2: int}
\infer1{\Gamma \vdash a_1 \upto a_2: set[int, min = a_1, max = a_2]}
% Is max/min enough here? what if we union two ranges with dijoint max/mins
\end{prooftree}
\\
\tag{Set Operations - Powerset}
\begin{prooftree}
\hypo{\Gamma \vdash a:T}
\infer1{\Gamma \vdash \mathbb{P}(a):set[T]}
\end{prooftree}
\end{gather}

\subsubsection{Bag Operations}
\paragraph{Includes:} Bag union, bag intersection, etc.

Let $\otimes = \Set{\cup, \cap, +, -}$
\begin{gather}
\tag{Bag Operations}
\begin{prooftree}
\hypo{\Gamma \vdash a_1, a_2: bag[T]}
\infer1{\Gamma \vdash a_1 \otimes a_2: bag[T]}
\end{prooftree}
\\
\tag{Bag Operations - Image}
\begin{prooftree}
\hypo{\Gamma \vdash a: T_1 \rel T_2, b: bag[T_1]}
\infer1{\Gamma \vdash a[b]: bag[T_2] \quad a[b] = a^{-1} \circ b}
\end{prooftree}
\end{gather}

\subsubsection{Relation Operations}
\paragraph{Includes:} Relation overriding, composition, domain/range manipulation. See Section 5.1.3 for changes in relational subtypes after applying operations.

\begin{gather}
\tag{Relation Operations - Function Call}
\begin{prooftree}
\hypo{\Gamma \vdash a: T_1 \rel T_2, b:T_1}
\infer1{\Gamma \vdash  a(b):T_2}
\end{prooftree}
\\
\tag{Relation Operations - Image}
\begin{prooftree}
\hypo{\Gamma \vdash a: T_1 \rel T_2, b: set[T_1]}
\infer1{\Gamma \vdash a[b]: set[T_2]}
\end{prooftree}
\\
\tag{Relation Operations - Overriding}
\begin{prooftree}
\hypo{\Gamma \vdash a_1, a_2: T_1 \rel T_2}
\infer1{\Gamma \vdash a_1 \oplus a_2: T_1 \rel T_2}
\end{prooftree}
\\
\tag{Relation Operations - Composition}
\begin{prooftree}
\hypo{\Gamma \vdash a_1: T_1 \rel T_2, a_2: T_2 \rel T_3}
\infer1{\Gamma \vdash a_1 \circ a_2: T_1 \rel T_3}
\end{prooftree}
\\
\tag{Relation Operations - Domain Restriction}
\begin{prooftree}
\hypo{\Gamma \vdash a_1: set[T_1], a_2: T_1 \rel T_2}
\infer1{\Gamma \vdash a_1 \domres a_2: T_1 \rel T_2}
\end{prooftree}
\\
\tag{Relation Operations - Domain Subtraction}
\begin{prooftree}
\hypo{\Gamma \vdash a_1: set[T_1], a_2: T_1 \rel T_2}
\infer1{\Gamma \vdash a_1 \domsub a_2: T_1 \rel T_2}
\end{prooftree}
\\
\tag{Relation Operations - Range Restriction}
\begin{prooftree}
\hypo{\Gamma \vdash a_1: T_1 \rel T_2, a_2: set[T_2]}
\infer1{\Gamma \vdash a_1 \ranres a_2: T_1 \rel T_2}
\end{prooftree}
\\
\tag{Relation Operations - Range Subtraction}
\begin{prooftree}
\hypo{\Gamma \vdash a_1: T_1 \rel T_2, a_2: set[T_2]}
\infer1{\Gamma \vdash a_1 \ransub a_2: T_1 \rel T_2}
\end{prooftree}
\end{gather}

\subsubsection{Sequence Operations}
\paragraph{Includes:} Concatenation

\begin{gather}
\tag{Sequence Operations - Concatenation}
\begin{prooftree}
\hypo{\Gamma \vdash a_1, a_2: sequence[T]}
\infer1{\Gamma \vdash a_1 \concat a_2: sequence[T]}
\end{prooftree}
\end{gather}

\subsubsection{Numerical Operations}
\paragraph{Includes:} Addition, subtraction, multiplication, etc. on natrual numbers, integers, and floats.

\begin{gather}
\tag{Integer Operations - Division}
\begin{prooftree}
\hypo{\Gamma \vdash a_1: T_1, a_2: T_2}
\hypo{T_i \in \Set{int, nat}}
\infer2{\Gamma \vdash a_1 \text{ div } a_2: max(T_1, T_2)}
\end{prooftree}
\\
\tag{Integer Operations - Modulo}
\begin{prooftree}
\hypo{\Gamma \vdash a_1: T_1, a_2: T_2}
\hypo{T_i \in \Set{int, nat}}
\infer2{\Gamma \vdash a_1 \text{ mod } a_2: max(T_1, T_2)}
\end{prooftree}
\\
\tag{Numerical Operations - Addition}
\begin{prooftree}
\hypo{\Gamma \vdash a_1: T_1, a_2: T_2}
\hypo{T_i \in \Set{int, nat, float}}
\infer2{\Gamma \vdash a_1 + a_2: max(T_1, T_2)}
\end{prooftree}
\\
\tag{Numerical Operations - Subtraction}
\begin{prooftree}
\hypo{\Gamma \vdash a_1: T_1, a_2: T_2}
\hypo{T_i \in \Set{int, nat, float}}
\infer2{\Gamma \vdash a_1 - a_2: max(T_1, T_2, int)}
\end{prooftree}
\\
\tag{Numerical Operations - Floating Division}
\begin{prooftree}
\hypo{\Gamma \vdash a_1: T_1, a_2: T_2}
\hypo{T_i \in \Set{int, nat, float}}
\infer2{\Gamma \vdash a_1 * a_2: float}
\end{prooftree}
\\
\tag{Numerical Operations - Multiplication}
\begin{prooftree}
\hypo{\Gamma \vdash a_1: T_1, a_2: T_2}
\hypo{T_i \in \Set{int, nat, float}}
\infer2{\Gamma \vdash a_1 \text{ div } a_2: max(T_1, T_2)}
\end{prooftree}
\\
\tag{Numerical Operations - Negation}
\begin{prooftree}
\hypo{\Gamma \vdash a_1: T}
\hypo{T \in \Set{int, nat, float}}
\infer2{\Gamma \vdash -a_1: max(int, T)}
\end{prooftree}
\\
\tag{Numerical Operations - Exponentiation}
\begin{prooftree}
\hypo{\Gamma \vdash a_1: T_1, a_2: T_2}
\hypo{T_i \in \Set{int, nat, float}}
\infer2{\Gamma \vdash a_1 \text{ \^{} } a_2: max(T_1, T_2)}
\end{prooftree}
\end{gather}

\subsubsection{Record Types}
\paragraph{Includes:} Types for named tuples.

\begin{gather}
\tag{Records - Access}
\begin{prooftree}
\hypo{\Gamma \vdash a: Record[I:A, ...]}
\infer1{\Gamma \vdash a.I: A}
\end{prooftree}
\\
\tag{Records - Type Definition}
\begin{prooftree}
\hypo{\Gamma \vdash A_1, ..., A_n}
\hypo{\Forall{I_i, I_j}{I_i \neq I_j}}
\infer2{\Gamma \vdash Record[I_1: A_1, ..., I_n: A_n]}
\end{prooftree}
\\
\tag{Records - Initialization}
\begin{prooftree}
\hypo{\Gamma \vdash a_1: A_1, ..., a_n: A_n}
\infer1{\Gamma \vdash record(a_1=A_1, ..., a_n=A_n): Record[a_1: A_1, ..., a_n:A_n]}
\end{prooftree}
\end{gather}

\subsubsection{Compound Commands}
\paragraph{Includes:} If, While, For, and Import statements, along with Procedure definitions and calls.

\begin{gather}
\tag{Command - Composition}
\begin{prooftree}
\hypo{\Gamma \vdash C_1, C_2}
\infer1{\Gamma \vdash C_1 \text{ \textbackslash n } C_2}
\end{prooftree}
\\
\tag{Command - Valid Import Module}
\begin{prooftree}
\hypo{\Gamma \vdash \diamond}
\infer1{\Gamma \vdash \text{import } M: C}
\end{prooftree}
\\
\tag{Command - Valid Import Names}
\begin{prooftree}
\hypo{\Gamma \vdash \diamond}
\infer1{\Gamma \vdash \text{from } M \text{ import } I_1, ..., I_n: C}
\end{prooftree}
\\
\tag{Command - Import Module}
\begin{prooftree}
\hypo{\Gamma \vdash \text{import } M \text{ (with identifiers and types } I_1: A_1, ..., I_n: A_n)}
\infer1{\Gamma, M=record(I_1=A_1, ..., I_n=A_n) \vdash \diamond}
\end{prooftree}
\\
\tag{Command - Import Names}
\begin{prooftree}
\hypo{\Gamma \vdash \text{from } M \text{ import } I_1, ..., I_n \text{ (with types } A_1, ..., A_n)}
\infer1{\Gamma, I_1:A_1, ..., I_n:A_n \vdash \diamond}
\end{prooftree}
\\
\tag{Command - Procedure Definition}
\begin{prooftree}
% Read this as: make a new scope, place the procedure I and params x_* to create the body C:R.
% Then, the current scope produces a symbol I with type procedure[...]
\hypo{\Gamma, I: procedure[A_1, ..., A_n, R], x_1: A_1, ..., x_n: A_n \vdash C: R}
\hypo{I, x_1, ..., x_n \notin dom(\Gamma)}
\hypo{i \neq j \implies x_i \neq x_j}
\infer3{\Gamma \vdash \text{ procedure }I=C : procedure[A_1, ..., A_n, R]}
\end{prooftree}
\\
\tag{Command - Procedure Call}
\begin{prooftree}
\hypo{\Gamma \vdash procedure[A_1, ..., A_n, R], a_1: A_1, ..., a_n:A_n}
\infer1{\Gamma \vdash I(a_1, ..., a_n): R}
\end{prooftree}
\\
\tag{Command - If}
\begin{prooftree}
% C is a new scope (block) but doesnt necessarily define new bindings
\hypo{\Gamma \vdash b: bool, \text{block }C}
\infer1{\Gamma \vdash \text{if } b: C}
\end{prooftree}
\\
\tag{Command - If Else}
\begin{prooftree}
\hypo{\Gamma \vdash b: bool, \text{block }C_1, \text{block }C_2}
\infer1{\Gamma \vdash \text{if } b: C_1 \text{ else}: C_2}
\end{prooftree}
\\
\tag{Command - While}
\begin{prooftree}
\hypo{\Gamma \vdash b: bool, \text{block }C}
\infer1{\Gamma \vdash \text{while } b: C}
\end{prooftree}
\\
\tag{Command - For}
\begin{prooftree}
\hypo{\Gamma \vdash G: set[T]}
\hypo{\Gamma, i: T \vdash C}
\hypo{i \notin dom(\Gamma)}
\infer3{\Gamma \vdash \text{for } i \in G: C}
\end{prooftree}
\\
\tag{Command - Block}
\begin{prooftree}
% Declarations D produce a new environment in C
\hypo{\Gamma \vdash D \therefore \Gamma'}
\hypo{\Gamma' \vdash C}
\infer2{\Gamma \vdash \text{block }D \text{ \textbackslash n } C}
\end{prooftree}
\end{gather}



\subsubsection{Built-in Functions}
\paragraph{Includes:} Common set-related functions.

\begin{gather}
\tag{Built-in - Minimum}
\begin{prooftree}
\hypo{\Gamma \vdash T[\text{with } order \in traits(T)]}
\infer1{\Gamma \vdash min: set[T] \rightarrow T}
\end{prooftree}
\\
\tag{Built-in - Mapped Minimum}
\begin{prooftree}
\hypo{\Gamma \vdash T, int}
\infer1{\Gamma \vdash min: set[T] \times (T \rightarrow int) \rightarrow T}
\end{prooftree}
\\
\tag{Built-in - Maximum}
\begin{prooftree}
\hypo{\Gamma \vdash T[\text{with } order \in traits(T)]}
\infer1{\Gamma \vdash max: set[T] \rightarrow T}
\end{prooftree}
\\
\tag{Built-in - Mapped Maximum}
\begin{prooftree}
\hypo{\Gamma \vdash T, int}
\infer1{\Gamma \vdash max: set[T] \times (T \rightarrow int) \rightarrow T}
\end{prooftree}
\\
\tag{Built-in - Choice}
\begin{prooftree}
\hypo{\Gamma \vdash T}
\infer1{\Gamma \vdash choice: set[T] \rightarrow T}
\end{prooftree}
\\
\tag{Built-in - Domain}
\begin{prooftree}
\hypo{\Gamma \vdash T}
\infer1{\Gamma \vdash dom: T_1 \rel T_2 \rightarrow set[T_1]}
\end{prooftree}
\\
\tag{Built-in - Range}
\begin{prooftree}
\hypo{\Gamma \vdash T}
\infer1{\Gamma \vdash ran: T_1 \rel T_2 \rightarrow set[T_2]}
\end{prooftree}
\\
\tag{Built-in - Cardinality}
\begin{prooftree}
\hypo{\Gamma \vdash T}
\infer1{\Gamma \vdash card: set[T] \rightarrow nat}
\end{prooftree}
\\
\tag{Built-in - Bag Size}
\begin{prooftree}
\hypo{\Gamma \vdash T}
\infer1{\Gamma \vdash size: bag[T] \rightarrow nat}
\end{prooftree}
\\
\tag{Built-in - Sum}
\begin{prooftree}
\hypo{\Gamma \vdash T}
\hypo{T \in \Set{int, float, nat}}
\infer2{\Gamma \vdash sum: set[T] \rightarrow T}
\end{prooftree}
\\
\tag{Built-in - Cast}
\begin{prooftree}
\hypo{\Gamma \vdash T_1, T_2}
\infer1{\Gamma \vdash cast: type[T_1] \times T_1 \rightarrow T_2}
\end{prooftree}
\\
\tag{Built-in - Cast With \footnote{$E$ represents `with' expressions to tag onto a type}}
\begin{prooftree}
\hypo{\Gamma \vdash T_1, T_2, x: E}
\infer1{\Gamma \vdash cast: type[T_1] \times T_1 \times set[E] \rightarrow T_2[\text{with } x: E]}
\end{prooftree}
\end{gather}
